\begin{proof}

\begin{enumerate}[label={$\langle1\rangle$\arabic*.},leftmargin=*]
\item By Ehresmann's theorem the smooth bundle over the contractible
disc is differentiably trivial.  Consequently
$H^1(\mathcal X,\mathbb Z)=H^2(\mathcal X,\mathbb Z)=0$.
The exponential sequence identifies
\[
 \operatorname{Pic}(\mathcal X)=H^1(\mathcal X,\mathcal O_{\mathcal X})
\]
as groups.  This is a complex vector space, so multiplication by
the nonzero integer $D$ is an automorphism.  This proves the
existence and uniqueness of the \emph{relative} root, including
when $D<0$.  The same argument on each fibre makes its Picard
group uniquely divisible.  \SeriesRef{push:src:E.2.1} gives
$K_{X_t}^{A}=(f_t^*\mathcal O(1))^{D}$, so uniqueness identifies
$\mathcal L_t$ with this actual pencil, not just with a numerical
class.

\item Proper direct-image coherence gives $\mathcal V=\pi_*\mathcal L$.
It is torsion-free: a local section killed by a nonzero function
of $t$ vanishes on the inverse image of an open set, hence
vanishes as a section of the invertible sheaf on the smooth
total space.  A coherent torsion-free sheaf on a smooth curve
is locally free.  Its rank is two because the noncentral
algebraic reductions have connected fibres.  Applying $\pi_*$
to
\[
 0\longrightarrow\mathcal L\xrightarrow{\,t\,}\mathcal L
 \longrightarrow i_*\mathcal L_0\longrightarrow0
\]
shows directly that $\mathcal V/t\mathcal V$ injects into
$H^0(X_0,\mathcal L_0)$.  Thus the two central sections are
linearly independent even if central cohomology jumps; no
surjectivity of base change has been assumed.

\item We use precisely the non-holomorphic-meromorphic-map theorem
of \citet[Theorem~2.1]{CDP2020D}, also written out in
\citet[Theorem~3.1 and Corollary~3.2]{LRS2019}.
For a compact complex threefold with $b_2=0$, a nonconstant
meromorphic map to a curve which fails to be holomorphic
would force $c_3=0$.  Since $c_3(X_0)[X_0]=2$, every
meromorphic function on $X_0$ extends as a holomorphic map
to $\mathbb P^1$.  This argument uses neither the different
holomorphic-fibration assertion of CDP, Theorem~2.2, nor its
Corollary~2.3.  Those assertions cannot be applied here.

\item In particular the ratio of $s_0,s_1$ gives a nonconstant
holomorphic map $h_0:X_0\to\mathbb P^1$.  Also $a(X_0)\le1$:
otherwise two algebraically independent meromorphic functions
would give a holomorphic map $X_0\to\mathbb P^1\times\mathbb P^1$
with two-dimensional image.  At a point of differential rank
two, after choosing finite affine target coordinates $u,v$
centred at its image, the global meromorphic function $u/v$
would have the local indeterminacy of the quotient of two
independent coordinates.  It cannot extend holomorphically
there.  This contradicts the preceding theorem.  Hence
$a(X_0)=1$.

\item Take the Stein factorization $h_0=R\circ g_0$.  Its intermediate
space is a smooth compact curve.  Connectedness of the fibres
injects its $H^1(\ ,\mathbb Z)$ into $H^1(X_0,\mathbb Z)=0$
by the degree-one Leray sequence, so it is $\mathbb P^1$.
Every meromorphic function on $X_0$ factors rationally through
$g_0$: the image of its holomorphic map together with $g_0$
is an analytic curve in $\mathbb P^1\times\mathbb P^1$;
normalizing this curve gives a finite factor of $g_0$, whose
degree is one because $g_0$ has connected fibres.  Thus $g_0$
is the algebraic reduction, and $R$ has some degree $e\ge1$.

\item Finally, on a target chart where the first homogeneous
coordinate of $h_0$ is nonzero, $s_1=h_0s_0$ in a local
trivialization, and the common ideal is generated by $s_0$.
The analogous assertion holds on the other chart.  These local
generators define a nonzero section of
$\mathcal L_0\otimes h_0^*\mathcal O(-1)$ and an effective
Cartier zero divisor $E_0$.  This gives the stated factorization
of $\mathcal L_0$ with no residual codimension-two base locus.

\end{enumerate}
\end{proof}
